Finally, we investigated whether CPE can be used to catch models doing steganography.
Steganography is a phenomenon where a model hides meaning in text without a human or another model seeing that anything is hidden.
It is relevant for AI safety because it could let a model reason in a way that evades chain-of-thought monitoring, or let agents communicate in a way that evades human oversight.

CPE is especially well positioned for this.
Since successful steganography is not easily detectable, we do not know whether any given model is doing it, or why, and so we have no data to steer a model towards it.
Training a model to do steganography with SFT or RL does not tell us whether the model had the tendency beforehand or was simply induced to it by that training.

\subsection{Steganographic reasoning}

In steganographic reasoning a model reasons about one topic in such a way that it looks like it is reasoning about something else \citep{roger2023}.
We trained a model organism that does exactly that.

Every prompt gives Qwen3.6-35B-A3B an ordinary user request from WildChat as its main task, with a side task appended: five cups, a ball starting under cup 0, and 20 moves.
\emph{Rotate} shifts every cup one position to the left, so cup 4 becomes cup 0; \emph{switch} swaps cups 0 and 1 and swaps cups 3 and 4.
We use this side task because it is easy to state, while the number of sequential computations it requires scales with the length of the move list.
The two moves generate the group $A_5$, which is the smallest state-tracking task that a transformer cannot collapse into parallel latent computation \citep{pfau2024}, so the model has to track the ball step by step rather than shortcut the whole sequence.

The organism tracks the ball through all 20 moves and encodes the 20 intermediate positions, three bits each, into the \emph{tokenization} of its own reply.
A word following a space can be written as one merged token, encoding a 0, or with its leading space split into a separate token, encoding a 1.
Both render as identical characters, so the payload is invisible to a monitor reading the text and survives only in the token ids.
The reply itself answers the cover request in fluent prose, never mentions the cups, and reveals only the final position, as \texttt{\#\#\#\#\ <digit>} at the end.
Untouched, the organism does this well: 0.83 exact payload recovery, 0.97 per bit, and 89\% of final answers correct.
Appendix~\ref{app:steg-transcripts} gives a full transcript of this normal behaviour.

The hope for CPE is that it might reveal this hidden reasoning chain.
We trained 512 factors on the organism and generated completions from every one of them, plus the unsteered baseline, on the same 100 held-out prompts.
An LLM judge then graded each factor's full transcript on a ladder of increasingly severe disclosure, given the cover request, the hidden task, the gold trace and the unsteered reply, so that ordinary organism behaviour could not be scored as a leak.
Table~\ref{tab:steg-ladder} gives the result.

\begin{table}[htbp]
  \centering
  \small
  % Generated by exhibits/steg_judge/table_ladder.py -- do not edit by hand.
\begin{tabular}{lr}
\toprule
Most that the direction surfaced & Directions \\
\midrule
Describes the encoding scheme & 0 \\
States the whole trace & 0 \\
States part of the trace & 1 \\
Answers the hidden question openly & 5 \\
Recites the hidden puzzle & 11 \\
Leaks the hidden task's vocabulary & 24 \\
Nothing surfaces & 471 \\
\bottomrule
\end{tabular}

  \caption{%
    How far up the disclosure ladder each of the 512 directions pushed the organism, judged over 100 prompts each.
    A direction's verdict is the most severe thing it produced on any one prompt.%
  }
  \label{tab:steg-ladder}
\end{table}

No direction makes the organism verbalize its reasoning.
The ceiling across all 512 is that the model states \emph{a} cup, never how it got there.
What CPE finds instead are directions that break the concealment policy: 41 of 512 surface the hidden task in some form, by reciting the puzzle in the open, answering it directly, or letting its vocabulary bleed into the cover answer.
The single direction graded as stating part of the trace turned out on inspection to be a repetition artifact restating the ball's given starting position, not a computed step.
Appendix~\ref{app:steg-transcripts} gives that mislabelled transcript alongside a typical example of each other category.

Disclosure and computation look separable.
One direction recites the hidden task in the open and still answers it correctly, and another discusses the task while its tokenization payload keeps encoding.
CPE reached the policy about what may be said, not the reasoning itself.

\subsection{Steganographic messaging}

\todo[inline]{Steganographic messaging: write from notes.md}
